\section{Replication Topologies}
\label{sec:replication-topologies}

When Alice renames a task on her laptop, which replica accepts the edit, and how does the new title reach Bob's phone?
In the server-centred task board, the application sends the write to a database leader.
In the local-first task board, Alice's laptop can accept the edit while offline.

A \emph{replication topology} describes the paths that changes follow between replicas.
In single-leader replication, the leader sends accepted writes to its followers.
Independently writable replicas can exchange their changes directly or through a synchronization server.
A server can therefore occupy a central place in either topology, but only the database leader must accept each write.
\parencites[152, 168, 175]{designing-data-intensive-applications}[173]{local-first-software}

\subsection{Single-leader replication}

In single-leader replication, one replica accepts writes and propagates them to the other replicas, called \emph{followers}.
The followers apply changes in the leader's order and can serve read requests.
Figure~\ref{fig:single-leader-replication} shows the write and replication paths.
\parencite[181--182]{understanding-distributed-systems}

In the server-centred task board, Alice's request to rename a task reaches the application server, which writes to the database leader.
The leader propagates the accepted change to its followers.
If the application cannot reach an available leader, it cannot complete the write through this replication model.
\parencite[152--153, 168]{designing-data-intensive-applications}

\begin{illustration}[htbp]
	\centering
	\caption{Single-leader replication}
	\label{fig:single-leader-replication}
	\begingroup
	\setstretch{1}
	\begin{tikzpicture}[
	thesis,
	client/.style={thesis/participant, minimum width=2.8cm, minimum height=0.7cm},
	replica/.style={thesis/store, minimum width=3.1cm, minimum height=0.8cm},
	network/.style={thesis/network, ->}
]
	\node[client] (client) at (0,0) {Client};
	\node[replica] (leader) at (0,-1.65) {Leader replica};
	\node[replica] (follower-a) at (-2.6,-3.5) {Follower replica 1};
	\node[replica] (follower-b) at (2.6,-3.5) {Follower replica 2};

	\draw[network] (client.south) -- (leader.north)
		node[midway, right=0.15cm] {Write};
	\draw[network] (leader.south west) -- (follower-a.north)
		node[midway, left=0.15cm] {Replicate};
	\draw[network] (leader.south east) -- (follower-b.north)
		node[midway, right=0.15cm] {Replicate};
\end{tikzpicture}

	\endgroup
	\figuresource{Author's illustration adapted from \textcite[182, Figure~19.1]{understanding-distributed-systems}.}
\end{illustration}
\FloatBarrier

The acknowledgement policy affects when the write can be reported as complete.
With asynchronous replication, the leader does not wait for followers to acknowledge the change, so a follower may still return the previous task title.
Synchronous replication waits for acknowledgements from the required followers, adding a dependency on their availability and response time.
\parencite[182--183]{understanding-distributed-systems}

\newpage
\subsection{Independently writable replicas}

Local-first editing requires devices to accept supported writes without reaching a remote leader.
\textcite[170--171]{designing-data-intensive-applications} describes offline clients and collaborative editing as uses of asynchronous \emph{multi-leader replication}: each device accepts local changes and exchanges them with other replicas later.
Accepting a local edit does not require coordination with other replicas before it completes.
However, replicas still need shared rules for interpreting and reconciling the changes they later exchange.

Suppose Alice's laptop and Bob's phone both hold a task titled \enquote{Draft}.
While disconnected, Alice renames it to \enquote{Review}; independently, Bob renames it to \enquote{Ready}.
Both changes are saved locally before either device learns about the other edit.
The replicas must subsequently exchange these changes, but independent writing does not require a direct connection between every pair of devices.

In a \emph{server-mediated star}, each client connects to a common synchronization server, as shown in Figure~\ref{fig:server-mediated-synchronization}.
Alice can send her change to the server, which retains data for Bob to retrieve when his phone reconnects.
This supports collaboration even when the devices are never online together.
The local-first essay calls such supporting servers \enquote{cloud peers}.
Their central position in the communication graph does not require clients to obtain approval for each local edit.
If the server becomes unavailable, exchange through it is delayed while the clients can continue editing their local data.
\parencite[173]{local-first-software}

\begin{illustration}[htbp]
	\centering
	\caption{Server-mediated synchronization between independently writable replicas}
	\label{fig:server-mediated-synchronization}
	\begingroup
	\setstretch{1}
	\begin{tikzpicture}[
	thesis,
	replica/.style={thesis/store, minimum width=3.8cm, minimum height=0.75cm},
	network/.style={thesis/network, <->}
]
	\draw[thesis/boundary] (0,0) rectangle (4.4,-1.5);
	\draw[thesis/boundary] (6.4,0) rectangle (10.8,-1.5);
	\node at (2.2,-0.3) {Alice's laptop};
	\node at (8.6,-0.3) {Bob's phone};
	\node[replica] (alice) at (2.2,-0.95) {Task-board replica\\Local reads and writes};
	\node[replica] (bob) at (8.6,-0.95) {Task-board replica\\Local reads and writes};

	\node[replica, minimum width=4.2cm] (server) at (5.4,-2.85)
		{Synchronization server\\Stored copy};
	\draw[network] (alice.south) -- (server.north west)
		node[pos=0.65, left=0.15cm] {Exchange changes};
	\draw[network] (bob.south) -- (server.north east)
		node[pos=0.65, right=0.15cm] {Exchange changes};
\end{tikzpicture}

	\endgroup
	\figuresource{Author's illustration based on \textcite[173]{local-first-software}.}
\end{illustration}
\FloatBarrier

With \emph{peer-to-peer communication}, devices exchange changes directly when a connection is available.
Alice and Bob could synchronize over a local network without routing their changes through a remote storage server.
However, the devices must discover one another and establish a connection; network configuration can make this difficult.
Direct exchange also requires both peers to be reachable at the same time.
A supporting server can retain data between connections, so direct peer links and server-mediated synchronization can be combined.
\parencite[173]{local-first-software}

Whether Alice and Bob exchange changes directly or through a server, their replicas must reconcile the independently edited titles.
The communication topology determines the paths those changes follow, while the reconciliation rules determine the resulting state.
